\section{Comprensión del negocio}
\label{sec:negocio}

Proyecto: relación entre el tipo de cambio oficial (BCB), el mercado paralelo
(\emph{dólar blue}) y el mercado P2P de USDT en Bolivia.

\subsection{Objetivo del negocio y KPIs}

\subsubsection{Definición del problema}

En Bolivia conviven \textbf{dos sistemas de precios del dólar} que no siempre
coinciden:

\begin{itemize}
  \item el \textbf{tipo de cambio oficial}, administrado por el Banco Central de
        Bolivia (BCB), que rige para operaciones bancarias formales, y
  \item el \textbf{mercado paralelo o ``dólar blue''}, que se forma libremente en
        plataformas P2P de criptomonedas (Binance, Bybit, Bitget) y
        fintechs (Airtm, Takenos, El Dorado, Wallbit), cotizado en USDT.
\end{itemize}

La \textbf{brecha} entre ambos es el principal termómetro de escasez de divisas
y de presión cambiaria. El problema que aborda este proyecto es de
\textbf{previsión}: los actores del mercado (importadores, empresas y personas)
no pueden anticipar el nivel futuro del paralelo, y por lo tanto no pueden
planificar tesorería ni cobertura con antelación. Se identifican dos problemas
específicos:

\begin{enumerate}
  \item \textbf{Falta de anticipación} --- no existen estimaciones del nivel del
        paralelo ni de la brecha a horizontes cortos (1~hora y 1~día). Predecir el
        tipo de cambio fuera de muestra es un problema reconocidamente difícil,
        en el que los modelos estructurales suelen no superar al paseo aleatorio
        ; aun así, enfoques de aprendizaje automático
        han mostrado capacidad predictiva en mercados financieros.
  \item \textbf{Relación no cuantificada} --- se desconoce si las normas
        publicadas en la Gaceta Oficial (combustibles, controles de cambio,
        reservas, deuda, etc.) anticipan o acompañan movimientos de la brecha;
        sin esa evidencia no pueden usarse como señales de política.
\end{enumerate}

Adicionalmente, la información está fragmentada entre fuentes heterogéneas
(BCB, bancos, agregadores y plataformas P2P), lo que condiciona la recolección
de datos y la construcción de una serie única de trabajo.

Definir con precisión este problema determina \textbf{qué datos recolectar}
(series de precios por fuente con marca de tiempo y origen, volúmenes) y
\textbf{cómo interpretar los análisis} (toda métrica se lee
como distancia frente al precio oficial administrado, no como precio absoluto),
y fija el criterio de éxito: la capacidad de \textbf{anticipar} el paralelo y
la brecha, y de vincular sus movimientos con las publicaciones de política.

\subsubsection{Objetivos}

\begin{table}[!ht]
\centering\small
\begin{tabular}{@{}p{0.15\textwidth}p{0.75\textwidth}@{}}
\toprule
\# & Objetivo \\
\midrule
Objetivo 1 & Predecir el nivel del paralelo y la brecha oficial--paralelo a
      1~h y 1~día, con intervalos de predicción \\
Objetivo 2 & Cuantificar la relación entre las publicaciones de la Gaceta
      (canales temáticos) y la brecha, en ventanas de $\pm5$ días \\
\bottomrule
\end{tabular}
\caption{Objetivos de predicción.}
\label{tab:objetivos}
\end{table}

\subsection{Criterios de éxito}

\begin{itemize}
  \item Los pronósticos a 1~h y 1~día superan a la línea base naïve (paseo
        aleatorio) fuera de muestra, medidos con RMSE y MAE
        \cite{makridakis2020m4}.
  \item Los intervalos de predicción alcanzan una cobertura empírica cercana
        al 90\%.
  \item El estudio de eventos de $\pm5$ días muestra un efecto medible (signo y
        magnitud, con errores HAC) de las publicaciones relevantes sobre la
        brecha, o una conclusión documentada de ausencia de relación.
\end{itemize}

\subsection{Inventario de recursos}

\subsubsection{Recursos --- datos disponibles}

El proyecto dispone de ocho conjuntos de datos: el historial completo de
snapshots (16,4 mil registros blue vs oficial), ~96 mil snapshots por
plataforma, la serie horaria alineada de 6 meses (4,3 mil puntos), historial
por plataforma cada 3~h, velas OHLC del P2P, percentiles y estadísticos
rodantes, el pipeline reproducible sobre 37 endpoints y el reporte de calidad
de la propia API.

\textbf{Herramientas y equipo:} Python 3.12 (\texttt{requests}, \texttt{pandas},
\texttt{matplotlib}, \texttt{seaborn}), análisis exploratorio ya avanzado (EDA
en PDF, auditoría de predictibilidad) y un analista con acceso a la API (clave
propia con cuota semanal).
